\section{Phase transitions with prescribed order parameters}

\begin{definition}[Discriminating systems, Georgii (16.18)--(16.20)]
    \label{def:op-16-18}
    \uses{def:erg-invariant, thm:erg-145}
    \lean{MeasureTheory.GibbsMeasure.withOne, MeasureTheory.GibbsMeasure.observableVector, MeasureTheory.GibbsMeasure.observableMatrix, MeasureTheory.GibbsMeasure.IsDiscriminating, MeasureTheory.GibbsMeasure.isDiscriminating_iff_span_eq_top, MeasureTheory.GibbsMeasure.isDiscriminating_extremePoints_invariantG, MeasureTheory.GibbsMeasure.le_encard_extremePoints_invariantG_of_isDiscriminating, MeasureTheory.GibbsMeasure.combo, MeasureTheory.GibbsMeasure.HasNondegenerateInvariantMeans, MeasureTheory.GibbsMeasure.invariantVariance, MeasureTheory.GibbsMeasure.sq_integral_lt_integral_sq_condExp_iff, MeasureTheory.GibbsMeasure.cond_mem_invariantG, MeasureTheory.GibbsMeasure.hasNondegenerateInvariantMeans_of_det_ne_zero, MeasureTheory.GibbsMeasure.integral_eq_of_condExp_invariantEvents_ae_eq_const, MeasureTheory.GibbsMeasure.isDiscriminating_invariantG_iff, MeasureTheory.GibbsMeasure.sum_smul_mem_invariantG, MeasureTheory.GibbsMeasure.sum_smul_mem_invariantFields, exists_linearIndependent_comp_of_span_eq_top, Matrix.det_ne_zero_iff_span_range_row_eq_top}
    \leanok

    With $f_0 = 1$, $\{f_1,\dots,f_N\} \subset \mathcal L$ is a discriminating system for a family
    $P$ of random fields if some $\mu_0,\dots,\mu_N \in P$ give a matrix $(\mu_m(f_n))$ of full rank,
    equivalently the vectors $(\mu(f_0),\dots,\mu(f_N))$, $\mu \in P$, span $\mathbb R^{N+1}$;
    $\Phi \ll \{f_1,\dots,f_N\}$ means $P = \mathcal G_\Theta(\Phi)$. (16.19): a discriminating
    system for $\mathcal G_\Theta(\Phi)$ discriminates $\operatorname{ex}\mathcal G_\Theta(\Phi)$
    (ergodic decomposition), so $|\operatorname{ex}\mathcal G_\Theta(\Phi)| \geq N+1$. Lemma
    (16.20): $\Phi \ll \{f_1,\dots,f_N\}$ iff some $\mu \in \mathcal G_\Theta(\Phi)$ satisfies
    $\mu(f_a)^2 < \mu(\mu(f_a\mid\mathcal I)^2)$ for all $a \neq 0$ (the positivity of a variance);
    proved for an arbitrary specification and any countable transformation group moving every
    finite volume off itself, (ii) $\Rightarrow$ (i) through the conditional measures
    $\mu(\cdot\mid A)$, $A \in \mathcal I$, rather than Georgii's finite subalgebras.
\end{definition}

\begin{theorem}[Georgii (16.21)--(16.25)]
    \label{thm:op-16-21}
    \uses{def:op-16-18, thm:gv-16-14, cor:gv-16-12, thm:tf-16-7, def:gv-16-10}
    \lean{MeasureTheory.GibbsMeasure.localPotentialB, MeasureTheory.GibbsMeasure.pairVolume, MeasureTheory.GibbsMeasure.pairObservable, MeasureTheory.GibbsMeasure.orderGenerators, MeasureTheory.GibbsMeasure.BThetaOrder, MeasureTheory.GibbsMeasure.isClosed_BThetaOrder, MeasureTheory.GibbsMeasure.correlation, MeasureTheory.GibbsMeasure.correlation_sub_sq_le, MeasureTheory.GibbsMeasure.tendsto_inv_card_sum_correlation, MeasureTheory.GibbsMeasure.hasNondegenerateInvariantMeans_of_specificEnergy_le, MeasureTheory.GibbsMeasure.exists_pos_lt_one_forall_lt_invariantVariance, MeasureTheory.GibbsMeasure.invariantVariance_eq_sum, MeasureTheory.GibbsMeasure.exists_mem_invariantG_specificEnergy_le, MeasureTheory.GibbsMeasure.specificEntropy_finsetSum_smul_ne_bot, MeasureTheory.GibbsMeasure.exists_pos_forall_exists_isDiscriminating, MeasureTheory.GibbsMeasure.exists_pos_forall_exists_isDiscriminating_BThetaOrder, MeasureTheory.GibbsMeasure.exists_isDiscriminating_of_isDiscriminating_invariantFieldsOfFiniteEntropy, MeasureTheory.GibbsMeasure.exists_isDiscriminating_and_le_encard_extremePoints, MeasureTheory.GibbsMeasure.exists_pos_forall_exists_isDiscriminating_norm_sub_lt, MeasureTheory.GibbsMeasure.exists_isDiscriminating_of_finite_excluded, MeasureTheory.GibbsMeasure.BThetaOrderNeg, MeasureTheory.GibbsMeasure.exists_pos_forall_exists_isDiscriminating_neg, MeasureTheory.GibbsMeasure.gasObservable, MeasureTheory.GibbsMeasure.isDiscriminating_gasObservable, MeasureTheory.GibbsMeasure.exists_isDiscriminating_gasObservable, MeasureTheory.GibbsMeasure.orderVolumes, MeasureTheory.GibbsMeasure.BThetaOrderCond, MeasureTheory.GibbsMeasure.isClosed_BThetaOrderCond, MeasureTheory.GibbsMeasure.eq_or_eq_neg_of_translate_pairVolume, MeasureTheory.GibbsMeasure.eq_or_eq_neg_of_forall_abs_inner_eq, Potential.restrictVolumes, MeasureTheory.GibbsMeasure.normAt_sub_restrictVolumes_le, MeasureTheory.GibbsMeasure.exists_mem_BThetaOrder_coe_eq_restrictVolumes, MeasureTheory.GibbsMeasure.BThetaOrderCond_eq_BThetaOrder, MeasureTheory.GibbsMeasure.mem_BThetaOrder_iff, Potential.IsPairPotential, MeasureTheory.GibbsMeasure.isPairPotential_of_mem_BThetaOrder, MeasureTheory.GibbsMeasure.isGasPotential_of_mem_BThetaOrder_gasObservable, MeasureTheory.GibbsMeasure.exists_isPairPotential_isGasPotential_isDiscriminating, MeasureTheory.GibbsMeasure.relativeEntropyIn_infinitePi, MeasureTheory.GibbsMeasure.specificEntropy_infinitePi, Potential.BTheta.invariantFieldsOfFiniteEntropy, Potential.BTheta.isBoundedBy_pressure_specificEnergyFunctional, Potential.BTheta.specificEntropy_toReal_le_negFenchel, Potential.BTheta.mem_invariantFieldsOfFiniteEntropy_of_mem_invariantG, Potential.localPotentialₗ, MeasureTheory.tendsto_inv_card_sum_integral_mul_vadd, lp.precomp, Filter.le_of_tendsto_inv_card_sum_of_forall_notMem}
    \leanok

    $\mathcal B_\Theta(f_1,\dots,f_N)$ is the closed linear span of $\Phi^{(\Lambda_0,f_n)}$ and
    $\Phi^{(\Lambda_i, f_m\,(f_n\circ\theta_i))}$, $\Lambda_i = \Lambda_0 \cup (\Lambda_0 - i)$,
    $i \in I$, Georgii's own ``more precise'' description, and this equals the subspace cut out
    by his three conditions on the interaction terms (his ``it is easily seen''): the latter is
    closed since norm convergence in $\mathcal B_\Theta$ forces uniform convergence of each term
    and the spans are finite-dimensional, the generators satisfy the conditions, and a potential
    satisfying them is the $\|\cdot\|_0$-limit of its truncations to finitely many translate
    classes, each a finite combination of generators; the step ``$\Lambda_i$, $\Lambda_j$ are not
    translates for $i \neq \pm j$'' needs only $\Lambda_0 \neq \emptyset$, not Georgii's cube, and
    $0 \notin I$ is his $i \neq 0$. Hence Comment (16.24)(3)'s first sentence (elements of
    $\mathcal B_\Theta(f)$ for single-site $f$ are pair potentials) and the gas-potential half of
    Example (16.25)(1); the Bernoulli field $\alpha^S$ has specific entropy
    $-\mathcal H(\alpha\mid\lambda)$. (16.22)--(16.23): the nondegeneracy of the invariant means
    passes from $\mu$ to any $\nu$ with $\langle\nu,\Psi\rangle \leq \langle\mu,\Psi\rangle +
    \varepsilon\|\Psi\|_0$ on $\mathcal B_\Theta(f)$ when $4\varepsilon|\Lambda_0|\|f_a\|^2$ stays
    below the variance of $\mu(f_a\mid\mathcal I)$, the mean ergodic correlation limit (14.A3)
    supplying $\lim |\Lambda|^{-1}\sum_i \mu(f_a\, f_a\circ\theta_i)$, and $\varepsilon$ is
    uniform in $a$ because the variance is a quadratic form. Theorem (16.21): for standard Borel
    $E$, $f_n \in \mathcal L_{\Lambda_0}$ and $\mu_0,\dots,\mu_N \in \mathcal P^\lambda_\Theta$
    of full rank, every $\Phi^0$ has a $\Phi \in \Phi^0 + \mathcal B_\Theta(f_1,\dots,f_N)$ with
    $\Phi \ll \{f_1,\dots,f_N\}$, hence $|\operatorname{ex}\mathcal G_\Theta(\Phi)| \geq N+1$
    (a); with $\|\Phi - \Phi^0\| < \delta$ once the $\mu_m$ are Gibbs for $\Phi^0$ (b). The core
    is stated for any closed convex cone containing the test potentials, which gives Comment
    (16.24)(1) (interaction terms on any finite family of volumes vanish, so the range is
    arbitrarily long) and (16.24)(2) (the ferromagnetic cone $\mathcal B^-_\Theta(f)$, $N = 1$).
    Example (16.25)(1): the lattice gas on $\{0,\dots,N\}$ with $f_n = 1_{\{\sigma_0 = n\}}$, the
    Dirac measures at the constant configurations giving full rank, and the modified potential a
    gas pair potential with vacuum $0$ as Georgii states. The criterion of (16.24)(3), (16.24)(4) and (16.25)(2) are
    Theorem \ref{thm:op-16-24}. Not formalised: the converse inclusion in (16.25)(1),
    (16.26)--(16.27).
\end{theorem}
\begin{proof}
    \leanok
\end{proof}

\begin{theorem}[Georgii (16.28)--(16.30), broken shift invariance]
    \label{thm:op-16-30}
    \uses{thm:op-16-21, thm:erg-14A3, thm:ext-77-concrete}
    \lean{MeasureTheory.GibbsMeasure.atInfinityWithin, MeasureTheory.GibbsMeasure.IsOscillating, MeasureTheory.GibbsMeasure.hasCountableBasis_atInfinityWithin, MeasureTheory.GibbsMeasure.isCountablyGenerated_atInfinityWithin, MeasureTheory.GibbsMeasure.exists_mem_extremePoints_invariantG_isOscillating, MeasureTheory.GibbsMeasure.correlation_sub_integral_sq_eq, MeasureTheory.GibbsMeasure.abs_correlation_sub_integral_sq_le, MeasureTheory.GibbsMeasure.integral_abs_condExp_tail_sub_integral_pos_of_isOscillating, MeasureTheory.GibbsMeasure.measure_condExp_tail_ne_integral_pos_of_isOscillating, MeasureTheory.GibbsMeasure.notMem_extremePoints_G_of_measure_condExp_tail_ne_integral_pos, MeasureTheory.GibbsMeasure.exists_mem_extremePoints_invariantG_notMem_extremePoints_G, MeasureTheory.GibbsMeasure.condExp_tail_ae_eq_integral_of_forall_condExp_comp_shift, MeasureTheory.GibbsMeasure.exists_mem_extremePoints_G_integral_comp_shift_ne, MeasureTheory.GibbsMeasure.exists_mem_extremePoints_G_notMem_invariantG_of_isOscillating, MeasureTheory.GibbsMeasure.nonempty_extremePoints_invariantG_diff_extremePoints_G, MeasureTheory.GibbsMeasure.nonempty_extremePoints_G_diff_invariantG, MeasureTheory.GibbsMeasure.exists_mem_invariantG_isOscillating, MeasureTheory.GibbsMeasure.altCfg, MeasureTheory.GibbsMeasure.altField, MeasureTheory.GibbsMeasure.spinObservable, MeasureTheory.GibbsMeasure.altAxis, MeasureTheory.GibbsMeasure.altField_mem_invariantFieldsOfFiniteEntropy, MeasureTheory.GibbsMeasure.correlation_spinObservable_altField, MeasureTheory.GibbsMeasure.isOscillating_altField, MeasureTheory.GibbsMeasure.exists_isOscillating_altAxis, MeasureTheory.GibbsMeasure.exists_nonempty_extremePoints_diff_altAxis, Filter.limsup_sub_liminf_le_of_eventually_abs_sub_le, Filter.liminf_lt_limsup_of_eventually_abs_sub_le}
    \leanok

    $\mu$ is $f$-$I$-oscillating if $\limsup_{|i|\to\infty, i\in I}\mu(f\,f\circ\theta_i) >
    \liminf$ (16.29), along the filter cofinite $\sqcap I$, which is countably generated. Lemma
    (16.28), for an arbitrary specification: (i) an $f$-$I$-oscillating
    $\mu \in \mathcal G_\Theta$ yields (ii) a $\nu \in \operatorname{ex}\mathcal G_\Theta$ with
    $\nu(\nu(f\mid\mathcal T) \neq \nu(f)) > 0$, hence $\nu \notin \operatorname{ex}\mathcal G$
    (a conclusion, not a hypothesis), by (14.17), (7.7) and L\'evy's downward theorem; (ii) yields
    (iii) a $\nu' \in \operatorname{ex}\mathcal G \setminus \mathcal G_\Theta$ with
    $\nu'(f\circ\theta_j) \neq \nu'(f)$, for bounded measurable $f$, by the mean ergodic theorem
    and the $(\mathcal G,\mathcal T)$-kernel (7.25). Theorem (16.30): for standard Borel $E$,
    $f \in \mathcal L_{\Lambda_0}$, $I$ infinite (symmetry unused) and an $f$-$I$-oscillating
    $\mu \in \mathcal P^\lambda_\Theta$, every $\Phi^0$ has a $\Phi \in \Phi^0 +
    \mathcal B^I_\Theta(f)$ with an $f$-$I$-oscillating shift-invariant Gibbs measure, so
    $\operatorname{ex}\mathcal G_\Theta(\Phi) \setminus \operatorname{ex}\mathcal G(\Phi)$ and
    $\operatorname{ex}\mathcal G(\Phi) \setminus \mathcal G_\Theta(\Phi)$ are nonempty. Erratum:
    the gap in (16.29) must exceed $4\varepsilon|\Lambda_0|\|f\|^2$, not $2\varepsilon|\Lambda_0|\|f\|^2$,
    since the limsup and the liminf each move by $2\varepsilon|\Lambda_0|\|f\|^2$; the statement
    is unaffected. Example (16.31): the field $(\delta_{\omega^+} + \delta_{\omega^-})/2$ with
    $\omega^\pm_j = \pm(-1)^{j_1}$ is shift invariant of finite specific entropy with
    $\mu(\sigma_0\sigma_i) = (-1)^{i_1}$, hence $\sigma_0$-$I$-oscillating along the first axis, so
    every $\Phi^0$ can be modified on the axis pairs into a potential with
    $\operatorname{ex}\mathcal G_\Theta \setminus \operatorname{ex}\mathcal G$ and
    $\operatorname{ex}\mathcal G \setminus \mathcal G_\Theta$ nonempty; Georgii's field is the
    product over lines of the alternating two-state chain, which agrees with this one on the axis,
    all that (16.30) uses.
\end{theorem}
\begin{proof}
    \leanok
\end{proof}

\begin{theorem}[Georgii Comment (16.24)(3)--(4) and Example (16.25)(2)]
    \label{thm:op-16-24}
    \uses{thm:op-16-21, def:op-16-18}
    \lean{InformationTheory.klDiv_cond_self, Submodule.span_eq_top_iff_forall_sum_mul_eq_zero, Submodule.span_ne_top_of_forall_sum_mul_eq_zero, Potential.BTheta.eq_of_isNormalized_of_isEquivalent, Potential.BTheta.not_isEquivalent_of_ne_of_isNormalized, MeasureTheory.GibbsMeasure.exists_integral_lt_integral_of_isDiscriminating, MeasureTheory.GibbsMeasure.exists_mem_invariantFieldsOfFiniteEntropy_hasNondegenerateInvariantMeans, MeasureTheory.GibbsMeasure.exists_isDiscriminating_neg_of_isDiscriminating_invariantFieldsOfFiniteEntropy, MeasureTheory.GibbsMeasure.precomp_eval_mem_localFunctionsOn, MeasureTheory.GibbsMeasure.specificEntropy_infinitePi_cond, MeasureTheory.GibbsMeasure.infinitePi_cond_mem_invariantFieldsOfFiniteEntropy, MeasureTheory.GibbsMeasure.absolutelyContinuous_map_eval_of_specificEntropy_ne_bot, MeasureTheory.GibbsMeasure.isDiscriminating_invariantFieldsOfFiniteEntropy_iff, MeasureTheory.GibbsMeasure.isDiscriminating_precomp_eval_iff, MeasureTheory.GibbsMeasure.exists_isDiscriminating_neg_integral_lt_integral, MeasureTheory.GibbsMeasure.avgOn_singleton_eq_integral_update, MeasureTheory.GibbsMeasure.isNormalized_of_mem_BThetaOrder, MeasureTheory.GibbsMeasure.coe_apply_eq_of_sub_mem_BThetaOrder, MeasureTheory.GibbsMeasure.isNormalized_of_sub_mem_BThetaOrder, MeasureTheory.GibbsMeasure.exists_isNormalized_isDiscriminating_norm_sub_lt, MeasureTheory.GibbsMeasure.not_isEquivalent_of_coe_apply_ne}
    \leanok

    Comment (16.24)(3), criterion: for single-site observables $f_n = f_n^*\circ\sigma_0$ the
    hypothesis of (16.21)(a) --- that $f_1,\dots,f_N$ discriminate $\mathcal P^\lambda_\Theta$ ---
    holds iff $1, f_1^*, \dots, f_N^*$ are linearly independent modulo $\lambda$-null sets, i.e.
    $\lambda(f_a^* \neq c) > 0$ for all $a \neq 0$ and all $c$. ``If'': the conditioned Bernoulli
    fields $\lambda(\cdot\mid A)^S$ lie in $\mathcal P^\lambda_\Theta$ with specific entropy
    $\log\lambda(A)$ (from $\mathcal H(\lambda(\cdot\mid A)\mid\lambda) = -\log\lambda(A)$) and realise
    the expectations $\lambda(1_A f_n^*)/\lambda(A)$, so a covector annihilating them makes
    $b_0 + f^*_{b'}$ integrate to $0$ over every measurable set; Georgii's disjointness of the $A_m$
    is not needed, only the rank. ``Only if'': $\mathcal k(\mu) > -\infty$ forces
    $\mathcal H_{\{0\}}(\mu) < \infty$, hence $\sigma_0\mu \ll \lambda$, hence $\mu(f_a) = c$ for
    every $\mu \in \mathcal P^\lambda_\Theta$, so the covector $(-c, a) \neq 0$ annihilates every
    row. Neither $E$ standard Borel nor $\iota$ nonempty is needed. Example (16.25)(2): a bounded
    measurable spin $s$ which is not $\lambda$-a.s. constant (Georgii's non-degenerate $\lambda$ on
    an even bounded $E \subset \mathbb R$, $s$ the inclusion) satisfies the criterion, so every
    $\Phi^0$ can be modified inside the ferromagnetic cone $\mathcal B^-_\Theta(\sigma_0)$, at
    arbitrarily long range, into a potential with two shift-invariant Gibbs measures of different
    magnetisation; the explicit description of $\mathcal B^-_\Theta(\sigma_0)$ as
    $\Phi_{\{i,j\}} = -J(i-j)\sigma_i\sigma_j$, $\Phi_{\{i\}} = -h\sigma_i$ is not proved, only the
    inclusion the conclusion uses. Comment (16.24)(4): if $\alpha(f_n^*) = 0$ then every potential
    of $\mathcal B_\Theta(f_1,\dots,f_N)$ is $\alpha$-normalized, since $\alpha_{\{j\}}$ is the
    average over the spin at $j$ and each generator carries an $\alpha$-mean-zero factor there;
    hence the potentials (16.21)(b) produces from $\alpha$-normalized $\Phi^0$ are
    $\alpha$-normalized, agree with $\Phi^0$ off the modified volumes, and distinct ones are
    non-equivalent by Theorem (2.35)(a) --- Georgii's infinite-dimensional family.
\end{theorem}
\begin{proof}
    \leanok
\end{proof}
